% Draft for section 3.1, based on A's classic implementation.
% Verify the final selected feature sets, cross-validation ranges and citations
% against A's frozen artifacts before incorporating it into the team paper.
\paragraph{Dictionary baseline (M0).}
We construct a dictionary from entity token sequences in the training data,
retaining a sequence only if all its observed entity occurrences share one
type. At inference time, exact token matches are selected greedily in
descending span length without overlap. This simple model measures the
performance obtainable from memorized training entities.

\paragraph{Token logistic regression (M1).}
We extract sparse handcrafted features for each token, including lexical,
orthographic, affix and contextual information, and predict BIO labels with a
multinomial logistic classifier. Features are fitted on training text and
predictions are made independently at each token. Invalid inside transitions
are repaired before extracting spans. This model tests whether features from
the surrounding text improve upon exact dictionary matching.

\paragraph{Linear-chain CRF (M2).}
A linear-chain conditional random field uses the shared handcrafted feature
pipeline while jointly scoring a sentence's label sequence and label
transitions. We use training-only, contiguous-block cross-validation for the
reported CRF model selection, then fit on the full training split. Unlike
independent token predictions, sequence modeling can learn regularities among
neighboring labels. Both M1 and M2 repair invalid BIO transitions before
entity scoring; counts of raw invalid transitions should therefore be reported
separately if used to assess sequence consistency.

% M1 and M2 settings and features must be checked against the selected A
% configuration for each domain. Adjacent model differences are descriptive,
% not causal estimates: M2, M3 and M4 differ in more than one design choice.
